\section{Related Work}

\noindent \textbf{Adversarial Attacks and Defenses on AI Models.}
Existing security research on AI models has predominantly focused on threats that manipulate the model's inference outcomes or training integrity. These include adversarial attacks~\cite{}, where imperceptible perturbations cause misclassification; data poisoning~\cite{}, which corrupts the training data to degrade performance; and backdoor attacks~\cite{}, which inject hidden triggers into models. While extensive defense mechanisms like adversarial training and robust distillation have been proposed for these threats, they fundamentally address the statistical vulnerabilities of machine learning algorithms. In contrast, our work focuses on the comparatively under-investigated domain of \textit{AI supply chain security}, where attackers exploit the model artifacts not to manipulate AI predictions, but to execute arbitrary, system-level malicious code on the victim's host machine.

\noindent \textbf{Malicious Code Poisoning in the AI Supply Chain.}
With the proliferation of open-source model hubs, the AI supply chain has become a lucrative target for attackers seeking to distribute malware by embedding it directly within model artifacts.
Early studies primarily focused on the vulnerabilities inherent in insecure serialization formats. Research has extensively documented how formats like Python's \texttt{pickle} can be exploited to execute arbitrary code during the deserialization process~\cite{}. To mitigate these threats, several static analysis and scanning tools have been developed. For instance, Fickling~\cite{} decompiles \texttt{pickle} bytecodes to identify suspicious patterns, while Picklescan~\cite{} and ModelScan~\cite{} integrate with antivirus engines like ClamAV to detect known malware signatures within model files. However, these defense mechanisms are predominantly static and rely heavily on pre-defined signatures or rules. Consequently, they are susceptible to evasion techniques, such as code obfuscation or the exploitation of previously unknown vulnerabilities, and often generate high false positive rates when legitimate model operations resemble malicious patterns.


\noindent \textbf{Framework-Native Attacks and API Abuse.}
Recent advancements in supply chain attacks have moved beyond simple serialization exploits towards more sophisticated, framework-native methodologies. These attacks leverage the legitimate, built-in APIs of AI frameworks to execute malicious behaviors, making them exceptionally stealthy.
Zhu et al.~\cite{} introduced TensorAbuse, a novel attack paradigm that exploits the hidden capabilities of TensorFlow APIs to construct persistent and stealthy attacks directly embedded within the model's computational graph. Because these attacks utilize officially supported framework operations rather than external system calls or standard malicious payloads, they effectively bypass traditional static security scanners. While the authors proposed an LLM-based tool to statically classify API capabilities and detect potentially exploitable APIs, this approach still operates at a static level and struggles with the semantic ambiguity of API usage, leading to potential inaccuracies in distinguishing between benign and malicious intents. In contrast, our proposed method, MAD, introduces a dynamic, cross-layer approach that monitors actual runtime behavior and leverages high-level execution context to accurately identify malicious intent, addressing the limitations of existing static analysis and signature-based defenses.

\noindent \textbf{System Call-Based Detection.}
Although largely unexplored in the context of AI models, system call sequence analysis is a cornerstone of traditional malware detection. Moving beyond early heuristic methods, recent studies leverage advanced deep learning to capture complex API behaviors. For instance, API2Vec+~\cite{cui202Xapi2vec} utilizes graph-based embeddings (temporal and attribute graphs) to capture intra- and inter-process contextual behaviors, effectively resolving API interleaving issues. Similarly, API-MalDetect~\cite{maniriho202Xapimaldetect} employs CNNs and BiGRUs to extract deep semantic features from long API sequences, while CruParamer~\cite{chen202Xcruparamer} incorporates parameter sensitivity into API embeddings to enhance detection accuracy. Inspired by the proven efficacy of system call profiling in traditional security, MAD adapts this low-level monitoring strategy to the AI domain, providing a robust, framework-agnostic foundation that cannot be easily evaded by application-layer obfuscation.

\noindent \textbf{Unsupervised Anomaly Detection.}
Unsupervised learning is highly effective for identifying zero-day anomalies without relying on labeled attack data. In the security domain, reconstruction-based models like AutoEncoders are widely utilized to establish benign behavioral baselines. For example, Yang et al.~\cite{} successfully applied ensemble learning with AutoEncoders to detect anomalous network traffic, and ContraMTD~\cite{} uses contrastive learning for malicious traffic scoring. Similar reconstruction and clustering techniques ~\cite{} have been deployed for hardware trojan detection and computer vision anomalies~\cite{}. Leveraging these unsupervised principles, MAD employs a sequence-to-sequence AutoEncoder to statistically model the highly regular, benign execution templates of AI frameworks. This enables our system to flag structural deviations caused by stealthy model supply-chain attacks without requiring any prior knowledge of specific malware signatures.
